## TEMP holding note — RQ3 compartment-archetype classification, re-run on current data (2026-08-16)

**Context**: re-implements an earlier notebook's compartment-level growth-trajectory classification
(`notebooks/growth_curve_attribution/av2_local_growth_curve_grouped_importance.ipynb`, last run
2026-08-08, before the environmental Set2-5 restructuring) as a standalone script against the
CURRENT data pipeline — `models/growth_curve_attribution/rq3_compartment_archetype_check.py`. Same
rule shapes as the notebook, percentile bands recomputed fresh on today's data (not copied from the
old run), plus two additions the notebook didn't have: each representative compartment's own mean
yield class and Chapman-Richards shape parameters (k, p), and what fraction of the curve's eventual
asymptote is reached at that compartment's mean survey age.

**Run as**: `python -m models.growth_curve_attribution.rq3_compartment_archetype_check --cohort 4survey`

### What happened — 4 of 5 categories work fine; 1 has a real, specific, fixable reason for being empty

4survey, 231 compartments, 56,526 plots (`apply_disturbance_cleaning=True`, the default used
throughout this project's target construction):

| Pattern | n compartments |
|---|---:|
| moderate mismatch / mixed | 159 |
| trajectory-break outlier | 26 |
| close to yldc / stable | 20 |
| consistently above yldc | 20 |
| consistently below yldc | 6 |
| **possible reset / measurement issue** | **0** |

**Diagnosed directly, not guessed**: `build_plot_level_table(cohort, apply_disturbance_cleaning=True)`
already excludes every plot flagged `any_clearfell_like`/`any_measurement_inconsistent` from the
population before this classification ever runs (by the design of that function — see its own
docstring: those plots have "no real trajectory to fit a fixed-shape curve to"). The "possible
reset / measurement issue" category is looking for exactly those flags at the compartment level
(>=10% of a compartment's plots) — but by the time it runs, none of the plots it's looking for are
still in the population. **Confirmed by re-running with `apply_disturbance_cleaning=False`**: the
category immediately populates (1 compartment, cpmt with sufficient concentration of flagged
plots), and the other categories shift only slightly (trajectory-break outlier: 26->29,
consistently above: 20->20, close to yldc: 20->18, consistently below: 6->5) — consistent with
adding ~300 previously-excluded plots back into a 56,000+ population, not a different story.

**Why the fixed version still only finds 1 compartment, not many — a real data characteristic, not
a remaining bug**: disturbance flags are genuinely rare project-wide (254 `any_clearfell_like` +
67 `any_measurement_inconsistent`, out of 56,841 plots without cleaning applied — about 0.55% of
plots total), and the compartment-level rule requires a 10%+ CONCENTRATION of flagged plots within
one compartment, not just their presence anywhere. Rare + requires local concentration = a
genuinely demanding bar, expected to be sparse even with the pipeline fixed.

**Conclusion**: 4 of the 5 rule categories are not too strict for current data — they populate with
real, substantial examples (6-26 compartments each) using the exact same threshold logic as the
original notebook, just with fresh percentile bands. The 5th category's emptiness has a specific,
now-understood cause (an upstream filtering step removing its target population before it runs),
not a flaw in the classification rules themselves or a mismatch with forestry ecology. If this
category is wanted for the write-up, re-run with `apply_disturbance_cleaning=False` specifically
for this diagnostic (methodologically fine here, since this classification is exploratory/
diagnostic, not feeding anything the disturbance-exclusion decision was designed to protect).

### Representative examples per category (`frac_of_asymptote_at_mean_age` = (1-exp(-k*Age))^p at
the compartment's own mean k/p/Age — how far up the sigmoid this compartment's data actually sits)

| Pattern | cpmt | n_plots | mean local_y_max_diff (m) | mean yldc | mean k | mean p | frac of asymptote at mean age |
|---|---:|---:|---:|---:|---:|---:|---:|
| close to yldc / stable | 1012 | 36 | +3.21 | 20.0 | 0.033 | 1.807 | 0.780 |
| consistently above yldc | 1098 | 302 | +16.24 | 12.1 | 0.031 | 1.739 | 0.402 |
| consistently below yldc | 2069 | 3 | -33.65 | 22.0 | 0.026 | 1.522 | 0.691 |
| trajectory-break outlier | 1053 | 66 | -4.71 | 11.1 | 0.035 | 1.198 | 0.876 |

**The single-example lead above does NOT hold up as a strong pattern once checked across the full
groups (2026-08-16)** — group means for `frac_of_asymptote_at_mean_age`:

| Pattern | n | mean | SD |
|---|---:|---:|---:|
| close to yldc / stable | 20 | 0.644 | 0.124 |
| consistently above yldc | 20 | 0.584 | 0.153 |
| consistently below yldc | 6 | 0.710 | 0.182 |
| moderate mismatch / mixed | 159 | 0.593 | 0.138 |
| trajectory-break outlier | 26 | 0.670 | 0.161 |

There's a weak, directionally-sensible signal — `consistently above yldc` has the lowest group mean,
`consistently below yldc` the highest, consistent with "caught earlier in growth" partly explaining
the above/below split — but the SDs (0.12-0.18) are large enough that the ranges overlap heavily
across every group. The single compartment originally pulled as the "above yldc" example (cpmt
1098, 0.402) was an unusually extreme case within its own category, not representative of the
group's actual mean (0.584) — corrected here rather than left standing as a stronger finding than
it is.

### Environmental means by archetype (2026-08-16) — corroborates the boundary-proximity finding at compartment scale

Merged `elevation`, `topex`, `windward_topex`, `slope_degrees`, `gwa_wind_speed_50m`,
`CanopyCover`, and the three boundary-distance columns onto the plot table before compartment
aggregation (`merge_environmental_features()`, unchanged). A targeted set, not the full candidate
pool -- chosen because each is already established elsewhere in RQ3's work as real signal worth
checking against this new classification.

| Archetype | n | elevation | topex | windward_topex | slope | CanopyCover | dist_to_cpmt_boundary | dist_to_block_boundary | dist_to_forest_perimeter |
|---|---:|---:|---:|---:|---:|---:|---:|---:|---:|
| close to yldc / stable | 20 | 114.1 | 2.31 | 0.73 | 7.48 | 0.74 | 30.1 | 31.1 | 196.0 |
| consistently above yldc | 20 | 175.5 | 8.91 | 2.29 | 12.93 | 0.78 | 24.2 | 25.1 | 133.2 |
| consistently below yldc | 6 | 185.2 | 5.04 | -0.98 | 11.83 | 0.49 | 13.6 | 15.5 | 126.3 |
| trajectory-break outlier | 26 | 220.6 | 0.87 | 2.66 | 12.80 | 0.57 | 19.7 | 20.7 | 153.1 |
| moderate mismatch / mixed | 159 | 179.1 | 2.48 | 0.81 | 10.02 | 0.71 | 28.0 | 29.3 | 174.7 |
| *(4survey population median)* | -- | 149.9 | 0.35 | 1.07 | 8.81 | 0.79 | 28.4 | 29.9 | 165.5 |

**Two real, corroborating patterns, checked against population medians, not eyeballed against
nothing**:
1. **`consistently below yldc` and `trajectory-break outlier` compartments sit notably closer to
   compartment/block boundaries than `close to yldc/stable`** (13.6-19.7m vs. 30.1m dist_to_cpmt_
   boundary, population median 28.4m) — extends the individual-outlier-plot boundary-proximity
   finding (originally found on the 10 hand-picked plots) to whole compartment archetypes, a
   different unit of analysis, independently.
2. **`trajectory-break outlier` compartments have the highest elevation (220.6m) and highest
   windward-specific exposure (`windward_topex`=2.66) of any group; `consistently above yldc` has
   the highest omnidirectional exposure (`topex`=8.91) instead.** A real, meaningful distinction
   between the two wind metrics, not the same signal twice — omnidirectional exposure associates
   with growing faster than the yldc benchmark predicts; windward-specific exposure associates more
   with unstable/break-type trajectories.

**Honest caveat**: these are directional patterns from group means (n=6 to 159 compartments), not a
significance-tested result — no test run to confirm they exceed what group-size-driven noise alone
would produce. Consistent in direction with, and a genuine corroboration of, findings established
elsewhere in this project, not proof on its own.

### Not yet done

The representative-compartment trajectory plots (observed height by age vs. mean fitted local curve
vs. mean official yldc curve, one panel per category) that the original notebook built — data is
all present in the merged plot table, plotting code not yet written for the new script.
